% ----- Consignes exo 8 ----- %
\begin{td-exo}[\textsc{Vertex Cover} avec un arbre de recherche en profondeur]\,\\ % 8 
	Cet exercice est basé sur l'article de Carla Savage.

	Nous utilisons un parcours en profondeur d'un graphe \(G\) à partir d'un sommet quelconque \(r\), soit \(T\) l'arbre obtenu.
	Le vertex cover retenu est l'ensemble des sommets internes de \(T\) (c'est à dire les sommets de \(T\) qui ne sont pas des feuilles).
	Dans le cas où \(r\) est une feuille on prend ce sommet.

	Dans la suite nous notera \(M(T)\) (respectivement \(M(G)\)) un couplage maximum dans \(T\) (respectivement \(G\)).
	Soit \(VC^\ast(G)\) un vertex cover optimal dans \(G\).

	\begin{enumerate}
		\item Expliquer pourquoi l'algorithme précédent donne une solution réalisable.

		\item Rappeler pourquoi nous avons \(M(T) \leq M(G) \leq VC^\ast(G)\).

		\item Soit \(L(T)\) la liste des feuilles de l'arbre \(T\). 
		Montrer par récurrence sur la profondeur de l'arbre que \(|L(T)| = \sum_{v \in NL(T)} (\text{fils}(v) - 1) + 1\) où \(NL(T)\) est la liste des sommets internes et \(\text{fils}(v)\) le nombre de fils du sommet \(v\) dans \(T\).
		Un sommet \(v\) est dit \defemph{free} par rapport à un couplage \(M\) s'il n'est pas couvert par \(M\).
		Soit \(X\) l'ensemble des sommets de \(T\) qui sont free par rapport à \(M(T)\).
		Montrer que \(|L(T)| \geq (|X| - 1) + 1 = |X|\).
		Conclure en montrant que 
		\begin{equation*}
			\frac{|NL(T)|}{|VC^\ast(G)|} \leq 2.
		\end{equation*}

		\item Est-ce que cette borne est atteinte?

		\item Montrer que si \(VC^\ast(G)\) est un vertex cover de \(G\) et \(T\) un arbre couvrant de \(G\) obtenu par parcours en profondeurs admettant \(VC^\ast(T)\) alors
		\begin{equation*}
			\frac{|VC^\ast(G)|}{|VC^\ast(T)|} \leq 2.
		\end{equation*}

		\item Montrer que la borne supérieure est atteinte.
	\end{enumerate}
\end{td-exo}

% ----- Solutions exo 8 ----- %
\iftoggle{showsolutions}{ 
	\begin{td-sol}[]\ % 8
		
	\end{td-sol}
}{}
